
\textbf{No model here has a predictive distribution.} These are point forecasts --- one number per
cell-month, no draws, no quantiles, no intervals. All nine are configured as deterministic
predictors, so a different export would not recover a distribution; it was never produced. Nothing
that scores a \textit{single model's} uncertainty --- CRPS of its predictive density, interval
coverage, density calibration --- can be computed from these files.

One qualification, since it would otherwise look like a contradiction in the output: the ensemble
tool can still produce a CRPS. It does so by treating the member models' point forecasts as an
empirical sample and scoring the spread across models. That is a property of the ensemble, not of
any model in it, and it is not the same quantity as a CRPS computed from a model's own posterior.
Worth knowing before the two are compared.

\textbf{Nothing here has been evaluated.} I checked only that the files are not broken: no
negatives, no missing values, no duplicate cell-months, the expected row counts, and totals within
the same order of magnitude as observed. That rules out a failed back-transformation or a collapsed
model and establishes nothing else.

\textbf{It is nine models, not eleven.} Two others were set up to draw 100 samples per forecast,
making them probabilistic counterparts of \texttt{dancing\_monkey} and \texttt{dark\_necessities}.
A point delivery needs one value per cell, and at one sample their configurations are
byte-identical to those siblings --- the same model twice. Running them as intended is blocked by
memory: 100 samples for 64,818 cells across thirteen origins needs far more RAM than any machine we
have.

\textbf{One observation will dominate any squared-error metric.} December 2021, cell 148759 has
\textbf{113,395} observed fatalities; the next largest anywhere in the window is 5,395. That month
falls in only one of the thirteen origins (\texttt{\dots\_12}), where that single cell is about
41\% of observed fatalities. I do not know whether it is a real event or a data error. Either way
it is better settled before the numbers are produced than explained afterwards.

\textbf{Column naming matches the tool.} These files use \texttt{pred\_lr\_ged\_sb}, which is the
ensemble tool's own default \texttt{target\_column}. The September delivery used
\texttt{pred\_lr\_sb\_best} instead, so if anything needs renaming when the two are compared it is
that one, not this. Renaming either takes minutes and needs no re-run.

\textbf{Comparison with the September delivery.} You were sent three variants per model then --- a
mean, a 95th percentile and a conditional mean --- each computed from that model family's
predictive distribution. These nine have no distribution, so only the mean has an equivalent here.
If the evaluation work has been built on one of the other two, the two deliveries are not
like-for-like. A correction while I am here: my September note said the full distributions had been
retained so alternative summaries could be produced without re-running. They were not kept.
